%%%%%%%%%%%%%%%%%%%%%%%%%%%%%%%%%%%%%
\chapter*{Lebenslauf}
%%%%%%%%%%%%%%%%%%%%%%%%%%%%%%%%%%%%%
%
%% No entry to the table of contents
%\addcontentsline{toc}{chapter}{Curriculum Vitae}
%
%% However, set a PDF bookmark:
\pdfbookmark[0]{Lebenslauf}{cv-de}
%
%Add a header title (Kolumnentitel setzen)
\markboth{Lebenslauf}{Lebenslauf}
%
%%%%%%%%%%%%%%%%%%%%%%%%%%%%%%%%%%%%%
%
%
%%%%%%%%%%%%%%%%%%%%%%%%%%%%%%%%%%%%%%%%%%%%%%%%%%%%%%%%%%%%
\todo[inline,caption={Deutschsprachigen Lebenslauf einfügen oder diese Seite ausblenden!}]{%
Beim Vorlegen der Dissertation an der KIT-Fakultät für Informatik muss ein Lebenslauf beigefügt werden.
Wird diese Seite nicht benötigt (\zB wenn man eine eigene Vorlage verwendet
oder wenn die Dissertation veröffentlicht wird), lässt sich diese Seite
durch die Umstellung des Schalters \printkeyword{showGermanCV}
auf \printkeyword{false} in der Datei \printfilepathquoted{./preambel/AlleSchalter.tex}
komplett ausblenden.
}%
%%%%%%%%%%%%%%%%%%%%%%%%%%%%%%%%%%%%%%%%%%%%%%%%%%%%%%%%%%%%
%
%
%%%%%%%%%%%%%%%%%%%%%%%%%%%%%%%%%%%%%%%%%%%%%%%%%%%%%%%%%%%%
\minisec{Persönliche Angaben}
%%%%%%%%%%%%%%%%%%%%%%%%%%%%%%%%%%%%%%%%%%%%%%%%%%%%%%%%%%%%
%% Durch Definition einer Gruppe mit \begingroup ... \endgroup 
%% sicherstellen, dass Anpassungen nur lokal wirken
%% und sich nicht auf die nachfolgenden tabellen erstrecken
\begingroup%
\small% Kleinere Schrift
\setlength{\tabcolsep}{0pt}% Kein Einzug
\renewcommand{\arraystretch}{1.0}% Abstand zwischen den Zeilen nicht vergrößern
\begin{longtable}{>{\raggedright\arraybackslash}p{0.4\linewidth-2\tabcolsep}>{\raggedright\arraybackslash}p{0.6\linewidth-2\tabcolsep}}%
%%-----------------------------------------------------------------------------------------------------
Titel: & \AuthorTitle \\
Name: & \AuthorName \\
Geburtsdatum: & 01.01.1990 \\
Geburtsort: & \PlaceOfBirth \\
Nationalität: & deutsch \\
%%-----------------------------------------------------------------------------------------------------
\end{longtable}
\endgroup%

%%%%%%%%%%%%%%%%%%%%%%%%%%%%%%%%%%%%%%%%%%%%%%%%%%%%%%%%%%%%
\minisec{Bildung}
%%%%%%%%%%%%%%%%%%%%%%%%%%%%%%%%%%%%%%%%%%%%%%%%%%%%%%%%%%%%
%% Durch Definition einer Gruppe mit \begingroup ... \endgroup 
%% sicherstellen, dass Anpassungen nur lokal wirken
%% und sich nicht auf die nachfolgenden tabellen erstrecken
\begingroup%
\small% Kleinere Schrift
\setlength{\tabcolsep}{0pt}% Kein Einzug
\renewcommand{\arraystretch}{1.3}% Etwas mehr Abstand zwischen den Zeilen, damit sie nicht zusammenfallen
\begin{longtable}{>{\raggedright\arraybackslash}p{0.4\linewidth-2\tabcolsep}>{\raggedright\arraybackslash}p{0.6\linewidth-2\tabcolsep}}%
%%-----------------------------------------------------------------------------------------------------
01.01.1900 -- 01.01.2000 & \textbf{Gymnasium XYZ in ABC}\\
%%-----------------------------------------------------------------------------------------------------
01.01.1900 -- 01.01.2000 & \textbf{Studium der XXX am KIT}\newline Abschluss als ...\\
%%-----------------------------------------------------------------------------------------------------
\end{longtable}
\endgroup%

%%%%%%%%%%%%%%%%%%%%%%%%%%%%%%%%%%%%%%%%%%%%%%%%%%%%%%%%%%%%
\minisec{Berufserfahrung}
%%%%%%%%%%%%%%%%%%%%%%%%%%%%%%%%%%%%%%%%%%%%%%%%%%%%%%%%%%%%
%% Durch Definition einer Gruppe mit \begingroup ... \endgroup 
%% sicherstellen, dass Anpassungen nur lokal wirken
%% und sich nicht auf die nachfolgenden tabellen erstrecken
\begingroup%
\small% Kleinere Schrift
\setlength{\tabcolsep}{0pt}% Kein Einzug
\renewcommand{\arraystretch}{1.3}% Etwas mehr Abstand zwischen den Zeilen, damit sie nicht zusammenfallen
\begin{longtable}{>{\raggedright\arraybackslash}p{0.4\linewidth-2\tabcolsep}>{\raggedright\arraybackslash}p{0.6\linewidth-2\tabcolsep}}%
%-----------------------------------------------------------------------------------------------------
01.01.1900 -- 01.01.2000 & \textbf{Stelle XXX}\newline Wissenschaftliche Hilfskraft\\
%-----------------------------------------------------------------------------------------------------
01.01.1900 -- 01.01.2000 & \textbf{Stelle YYY}\newline Berufsbezeichnung\\
%%-----------------------------------------------------------------------------------------------------
\end{longtable}
\endgroup%
